\section{Implementation Details}

\subsection{The Worker Implementation}
The Worker component is designed using the \textbf{Strategy Pattern}. This ensures loosely coupled architecture where the simulation logic is strictly separated from the communication logic.

\subsubsection{Simulator Interface}
The \texttt{Simulator} interface defines a contract for all physics models:
\begin{verbatim}
type Simulator interface {
    Run(seed int64, iterations int64) ([]float64, error)
}
\end{verbatim}
This allows us to seamlessly swap the \texttt{RandomWalk3D} implementation for more complex models (e.g., Ising Model or Lattice QCD) without modifying the gRPC client code.

\subsubsection{Random Walk Logic (The Physics)}
We simulate a particle undergoing Brownian motion in 3D space. The position $P_n$ at step $n$ is given by:
\[
P_n = P_{n-1} + \vec{\delta}
\]
Where $\vec{\delta}$ is a random vector where each component $\delta_x, \delta_y, \delta_z \in [-0.5, 0.5]$.

To guarantee \textbf{scientific reproducibility}, the Random Number Generator (RNG) is initialized with a specific \texttt{seed} transmitted by the Master.
\[
RNG_{state} = f(seed)
\]
\[
RNG_{state} = f(seed)
\]
This ensures that re-running Task $T_i$ always yields the exact same coordinate path.

\subsection{The Scheduler Implementation}
The Scheduler acts as the central coordinator and is responsible for ensuring data integrity and system resilience.

\subsubsection{Concurrency Control (Mutex)}
Since multiple workers request tasks simultaneously, the internal task queue is a \textbf{Critical Section}. To prevent Race Conditions (e.g., assigning the same task to two workers), we use a \texttt{sync.RWMutex}.
\begin{itemize}
    \item \textbf{Lock}: Applied when dequeuing a task or updating its status.
    \item \textbf{Unlock}: Released immediately after the memory operation to ensure high throughput.
\end{itemize}

\subsubsection{Fault Tolerance (The Reaper)}
The system implements a \textbf{Timeout-Based Recovery} mechanism. A background process (The Reaper) wakes up every 10 seconds to scan the state of all running tasks.
\[
\text{If } (T_{now} - T_{started}) > 30s \implies \text{Reset to QUEUED}
\]
State Diagram:
\[
\texttt{QUEUED} \xrightarrow{GetTask} \texttt{RUNNING} \xrightarrow{SubmitResult} \texttt{COMPLETED}
\]
\[
\texttt{RUNNING} \xrightarrow{Timeout} \texttt{QUEUED} \ (\text{Re-entry})
\]
\[
\texttt{RUNNING} \xrightarrow{Timeout} \texttt{QUEUED} \ (\text{Re-entry})
\]
This ensures that if a worker crashes or interacts with a "Black Hole" (network partition), the science is not lost.

\section{Deployment \& Observability}
To satisfy the ``Cloud-Native'' requirement, the system is fully containerized.

\subsection{Containerization}
We use \textbf{Multi-Stage Docker Builds} to minimize the attack surface and image size.
\begin{enumerate}
    \item \textbf{Build Stage}: Uses the \texttt{golang:1.24} image to compile the source code.
    \item \textbf{Runtime Stage}: Copies only the binary to a ``Distroless'' image (Google's stripped-down Linux). The final image is $<20$MB.
\end{enumerate}

\subsection{Orchestration}
\texttt{Docker Compose} manages the cluster lifecycle. It creates a private virtual network where:
\begin{itemize}
    \item The \texttt{Scheduler} listens on port 50051.
    \item 5 \texttt{Workers} automatically discover the scheduler via DNS (\texttt{scheduler:50051}).
    \item \texttt{Prometheus} scrapes metrics from the scheduler on port 2112.
\end{itemize}

\subsection{Metrics}
We expose a Prometheus-compatible \texttt{/metrics} endpoint. The key metric is \texttt{tasks\_completed\_total}, which provides a real-time view of the grid's throughput (Events Per Second).
